\input{estilo}
% La figura 8 del paper de MillWheel, con sus rótulos en inglés: el binario,
% el código del usuario, las tres APIs y el estado persistente.
\begin{document}
\begin{figura}\begin{tikzpicture}[
    metodo/.style={caja, fill=white, minimum width=30mm, minimum height=7mm, font=\ttfamily\small},
    api/.style={caja, minimum width=19mm, minimum height=11mm, align=center}]
  \draw[marco] (0,0) rectangle (11,6.6);
  \node[etiq, anchor=north east] at (10.85,6.5) {MillWheel system binary};
  \draw[marco, fill=white] (3.0,1.4) rectangle (10.2,5.6);
  \node[etiq, anchor=north east] at (10.05,5.5) {user code: computation};
  \node[metodo] (pr) at (3.2,4.5) {ProcessRecord};
  \node[metodo] (pt) at (3.2,3.4) {ProcessTimer};
  \foreach \n/\x/\t in {st/4.7/State, ti/6.9/Timer, pd/9.1/Produce}
    \node[api] (\n) at (\x,1.4) {\t\\API};
  \draw[datos] (-1.6,4.5) -- node[etiq, above] {RPCs} (pr.west);
  \draw[datos] (11,3.0) -- node[etiq, above] {RPCs} (12.6,3.0);
  \draw[fill=relleno] (3.6,-2.3) -- (3.6,-3.5) arc[start angle=180, end angle=360, x radius=3.3, y radius=0.3] -- (10.2,-2.3);
  \draw[fill=relleno!50!white] (6.9,-2.3) ellipse[x radius=3.3, y radius=0.3];
  \node[etiq] at (6.9,-3.1) {persistent state};
  \foreach \n in {st,ti,pd}
    \draw[{Stealth[length=3mm,width=3mm]}-{Stealth[length=3mm,width=3mm]}, line width=1.6pt] (\n.south) -- (\n.south |- 0,-2.45);
\end{tikzpicture}\end{figura}
\end{document}
